\subsection{Decomposing changes in equilibrium strategies}

The analysis compares four pairs of equilibria at the ordinary cost level:
American to British fee shifting with risk neutrality held fixed; American to
British fee shifting with moderate risk aversion held fixed; risk neutrality to
moderate risk aversion under the American rule; and the same preference change
under the British rule. Both parties' preferences change in the latter two
comparisons. The tables report strategies from the original computed equilibria,
not the more broadly mixed profiles used in the sensitivity analysis.

\paragraph{Information sets and measurement.}
Strategies are compared at information sets reached with positive probability
in both endpoint equilibria. Each row conditions on the acting party's own
signal and observed history. The signal is on the common plaintiff-case-strength
scale, so a higher signal is more favorable to the plaintiff and less favorable
to the defendant. Offer information sets additionally distinguish the party's
private commitment to continue or exit if bargaining fails. A commitment to
abandon or default is not an unconditional probability of actual abandonment or
default, because settlement can occur before the commitment takes effect.

Filing, answering and exit commitments are reported as probabilities, expressed
as percentages. Their changes and decomposed contributions are expressed in
percentage points (pp): a change from 87.8\% to 100\% is an increase of
12.2 percentage points, not a 12.2\% proportional increase. Pure demands and
offers, and their decomposed contributions, are expressed as fractions of the
normalized award. For example, a demand changing from 0.75 to 0.05 changes by
$-0.70$. Offer amounts are used only when the endpoint strategies and the
counterfactual strategies entering the accounting are pure. Mixed offers are
not replaced by mean offers in these tables.

\paragraph{Best-response comparisons.}
Let $\theta_0$ and $\theta_1$ denote the original and target fee-rule or preference
specifications, and let $\sigma^0$ and $\sigma^1$ denote their selected equilibrium
strategy profiles. For each player, the opponent's strategy is partitioned into
three components: entry (filing or answering), offers, and exit commitments.
Under $\theta_1$, all eight combinations of original and target opponent
components are constructed. Each comparison holds the opponent's complete
hybrid strategy fixed and computes an unrestricted best response for the focal
player, jointly reoptimizing that player's current and subsequent decisions.
The calculation does not impose monotonicity or simulate an observed
adjustment process.

Write $z(S)$ for the reported strategy coordinate in this best response when
the set $S$ of opponent components has been replaced by its target-equilibrium
values. The coordinate is either the probability of the named binary action or
the pure demand or offer amount. Write $x$ and $y$ for the corresponding original
and target equilibrium coordinates. The direct contribution is
\[
  D=z(\varnothing)-x.
\]
For each of the six orders of replacing the three opponent components, the
marginal change from replacing component $k$ is recorded. Its contribution is
the average of those increments:
\[
  C_k=\frac{1}{6}\sum_{\pi}
  \left[z\bigl(S_{\pi,k}\cup\{k\}\bigr)-z(S_{\pi,k})\right],
\]
where $S_{\pi,k}$ contains the components preceding $k$ in order $\pi$. Thus
\[
  y-x=D+C_{\mathrm{entry}}+C_{\mathrm{offers}}+C_{\mathrm{exit}}+R,
  \qquad R=y-z(\{\mathrm{entry},\mathrm{offers},\mathrm{exit}\}).
\]
This is a direct-first accounting convention. Interactions between the change
in rules or preferences and the opponent's adjustment enter the opponent
contributions. These quantities are not uniquely identified causal shares.
Individual contributions can exceed the net change or have the opposite sign
because other contributions offset them.

\paragraph{Selection of the displayed changes.}
The displayed rows are a restricted subset of actual strategy changes. A changed
information set qualifies for the focus criterion if, against the target
opponent and with subsequent own decisions optimized, the original local
strategy assigns probability mass greater than $10^{-6}$ to actions whose
conditional utility is more than $10^{-6}$ below the best action. The utility
threshold is in the target specification's reported utility units. This is a
conditional action comparison, not the loss from reverting the entire strategy
or a welfare comparison across preference specifications. In particular,
disjoint strategy supports do not suffice: an action absent from the target
support can nevertheless remain optimal.

The criterion is evaluated for the original equilibrium profiles and for two
tested sets of equilibrium-preserving mixed representatives. The publication
tables retain its intersection across these three comparisons: 39 information
sets across the four interventions, displayed in 26 grouped rows. The counts by
intervention are 2, 11, 20 and 6 information sets. Consecutive signals are grouped
when their original and target distributions and numerical contributions agree
within $10^{-6}$, even if their sensitivity flags differ. ``At some signals''
means that at least one, but not every, signal included in that grouped row has
a tie- or off-path-sensitive allocation.

The original-profile accounting has no remaining term greater than $10^{-6}$
in the reported coordinate for these selected rows. This is verified before
producing a layout without a Remaining column; the remainder is not allocated
to another component. Selection for this display does not imply that other
strategy changes are unimportant or explained. Changes involving newly reached
or no-longer-reached histories, and redistribution among actions that remain
optimal, are outside the restricted table.

\paragraph{Equilibrium-preserving mixing checks.}
The auxiliary mixing search increases the mean normalized quadratic mixing
score $(1-\sum_a p_a^2)/(1-1/A)$ over information sets reached in the original
equilibrium, where $A$ is the full number of available actions. A player's
same-decision information sets are varied jointly. Candidate actions comprise
current support and actions tied in conditional utility under current
continuation play. Each proposed complete profile is checked against both
players' unrestricted best responses; a utility tie for the acting player alone
does not authorize adding an action to equilibrium support. Source-unvisited
policies remain fixed.

The main search admits maximum unilateral root-payoff gain of $10^{-9}$ reported
utility units and evaluates forward and reverse decision-block orders. The
higher-scoring verified result is selected. A second calculation uses forward
order and a tenfold tighter gain limit; its conditional tie threshold is also
tightened from $10^{-10}$ to $10^{-11}$. The search is local and does not certify
a globally maximally mixed or unique equilibrium. The selected reverse-order
search for the moderately risk-averse British case reaches the six-sweep limit.
The intersection therefore establishes stability only across the tested
representations. Even within that intersection, the decomposition can change
when a different mixed representative is used.

\paragraph{Ties, reach and verification.}
Primary best responses retain and renormalize the original probability on
actions that remain optimal within the $10^{-10}$ tie tolerance; when no such
probability remains, the first optimal action is selected. Alternative low- and
high-action selections test tie sensitivity. Separate checks replace opponent
policies unvisited in their donor equilibrium with extreme low- or high-action
completions. These are sensitivity checks, not bounds over all admissible
completions. The sensitivity column reports whether either check changes a
contribution in the original-profile accounting. It does not certify stability
under all equilibrium-preserving mixtures.

An intermediate best response may not actually reach an information set that is
reached in both endpoint equilibria. Conditional continuation comparisons are
still defined when chance-and-opponent reach is positive; the focal player's
own prior action probabilities are omitted from that reach calculation. If
chance-and-opponent reach is zero, no conditional allocation is imputed. The
underlying diagnostics retain these reach distinctions. Complete best responses
are independently replayed to check root utilities and conditional action
values, and the saved equilibrium profiles and reports are checked against the
initialized game. Reported table entries are rounded; accounting identities are
verified before rounding.
